\chapter{Valutazione Euristica degli Esperti}
\label{cap:valutazione_euristica}

In questo capitolo viene affrontata la fase di ispezione dell'usabilità di DietiEstates25 condotta 
attraverso il giudizio di esperti. A differenza dei test con gli utenti, questo approccio metodologico 
prevede che specialisti del dominio o di Human-Computer Interaction analizzino l'interfaccia per 
diagnosticare potenziali violazioni dei principi di design e prevedere le difficoltà che l'utente 
finale potrebbe incontrare.

\vspace{0.3cm}

La tecnica ispettiva adottata è la \textbf{Valutazione Euristica}, applicando come framework di 
riferimento le celebri \textbf{"Eight Golden Rules of Interface Design" (Otto Regole d'Oro)} teorizzate 
da Ben Shneiderman. Questi principi costituiscono una guida fondamentale per garantire un'interazione 
efficace e priva di frustrazioni:

\begin{enumerate}
    \item \textbf{Perseguire la coerenza (Consistency):} Garantire uniformità nell'uso di terminologia, 
    icone, colori, layout e flussi operativi in ogni sezione dell'applicazione.

    \item \textbf{Garantire l'usabilità universale (Shortcuts):} Soddisfare sia i novizi che gli utenti 
    esperti, integrando acceleratori (come scorciatoie da tastiera o menu rapidi) che velocizzino 
    l'interazione per chi ha familiarità con il sistema.

    \item \textbf{Offrire feedback informativo:} Il sistema deve comunicare costantemente il proprio 
    stato all'utente, fornendo risposte visive o testuali immediate a ogni azione compiuta.

    \item \textbf{Progettare dialoghi chiusi (Closure):} Le sequenze di azioni devono avere un inizio, 
    una parte centrale e una fine ben definita, offrendo all'utente un senso di compimento al termine 
    di un task (es. messaggio di successo).

    \item \textbf{Prevenire e gestire gli errori:} L'interfaccia deve essere disegnata in modo da 
    impedire agli utenti di commettere errori gravi e, qualora accadano, il sistema deve offrire 
    meccanismi semplici per rilevarli e correggerli.

    \item \textbf{Garantire la reversibilità delle azioni (Undo):} Le operazioni non devono essere 
    definitive; l'utente deve poter tornare indietro o annullare un'azione ("via di uscita di emergenza") 
    senza timore di danneggiare il lavoro svolto.

    \item \textbf{Supportare il Locus of Control interno:} L'utente deve percepire di essere l'iniziatore 
    delle azioni e di avere il pieno controllo del sistema, evitando sorprese o cambiamenti di stato 
    non richiesti.

    \item \textbf{Ridurre il carico della memoria a breve termine:} L'interfaccia deve favorire il 
    riconoscimento piuttosto che il ricordo, limitando la quantità di informazioni che l'utente deve 
    tenere a mente durante la navigazione.
\end{enumerate}

\vspace{0.3cm}

Dall'analisi ispettiva condotta sulle nostre interfacce applicando questi criteri, sono emerse le 
seguenti considerazioni:
